% paper/sections/07d_cls_stages.tex
%
% §5.2  CLS の A→F 6 段階アルゴリズム (CHK-223 NEW: §5 CLS-specific 配下)
%        §6.1 (CLS ψ FCCD 移流) / §5.1 (再初期化) を統合し,
%        1 物理時間ステップ内での順序を完結させる．
%        Stage A の TVD-RK3 更新式は §7.3 (eq:cls_tvd_rk3_psi) に集約 (CHK-223 Phase 3).
%        velocity-PPE ordering 詳細は §7.3 (sec:cls_velocity_consistency_v7) に集約.
%
% ═══════════════════════════════════════════════════════════════════════════════
\subsection{\texorpdfstring{CLS の A$\to$F 6 段階アルゴリズム}{CLS の A→F 6 段階アルゴリズム}}
\label{sec:cls_stages}
% ═══════════════════════════════════════════════════════════════════════════════

\noindent\textbf{Part 1 連続定式化との関係：}
本節は §~\ref{sec:cls_advection}（§3.2 CLS 保存形移流方程式）と
§~\ref{sec:reinit}（§3.3 再初期化方程式）で確立した連続定式化を
1 物理時間ステップの演算子分裂に組み込む離散アルゴリズムを示す．
本稿では，Stage A の FCCD ψ 移流
（§~\ref{sec:advection_motivation}）と Stage C の Ridge-Eikonal 再初期化
（§~\ref{sec:xi_sdf_reinit}）が主経路となる．

本節では，CLS \cite{OlssonKreiss2005} を 1 物理時間ステップ内の
6 段階に分解し，
\textbf{質量保存・幾何（法線・曲率）・質量補正}の
3 責務を独立に処理することで
既存 CLS の実装ミスを原理的に回避する設計を示す．

% -------------------------------------------------------------------------------
\subsubsection{\texorpdfstring{3 責務原則と $\psi\leftrightarrow\phi$ 対応}{3 責務原則と ψ↔φ 対応}}
\label{sec:cls_three_responsibilities}
% -------------------------------------------------------------------------------

\begin{center}
\begin{tabular}{@{}lll@{}}
\hline
責務 & 場 & 機構 \\
\hline
質量保存 & $\psi\in[0,1]$ & 保存形面フラックス移流 \\
幾何（法線・曲率）& $\phi$（符号距離）& 再初期化後に $\bnabla,\bnabla\cdot$ \\
質量補正 & $\psi$ への補正 & $\psi(1-\psi)$ 界面局在 \\
\hline
\end{tabular}
\end{center}

3 責務の混同——非保存形で $\phi$ を移流する，再初期化なしで $\psi$ から曲率を計算する，
質量補正を $\phi$ に適用する——が CLS 実装失敗の主要モードである．

\noindent\textbf{$\psi\leftrightarrow\phi$ 対応}（tanh/artanh ペア）：
\begin{equation}
  \phi = 2\varepsilon\,\mathrm{artanh}(2\psi-1),
  \qquad
  \psi = \tfrac{1}{2}\bigl(1 + \tanh(\phi/2\varepsilon)\bigr).
  \label{eq:cls_psi_phi_map}
\end{equation}
artanh 特異性回避のため artanh 前に $\psi\in[\delta,1-\delta]$
（$\delta=10^{-6}$ 程度）でクランプ．
$\varepsilon$ はシミュレーション全体で固定（変更すると $\kappa$ が変わる）．

% -------------------------------------------------------------------------------
\subsubsection{6 段階シーケンス}
\label{sec:cls_stage_sequence}
% -------------------------------------------------------------------------------

1 物理時間ステップ（$t^n\to t^{n+1}$）の段階順：

\begin{center}
\begin{tabular}{@{}llp{0.55\linewidth}@{}}
\hline
Stage & 操作 & 要件 \\
\hline
A & TVD-RK3 保存形移流 & FCCD 面値演算子による保存形フラックス（Stage A; \S\ref{sec:advection_motivation}），発散フリー $\mathbf{u}^n$．BC：壁面 $\partial_n\psi=0$（Neumann）/ 周期境界（検証），CCD ghost-cell 対称延長 \\
B & 移流後 $\psi(1-\psi)$ 補正 & 任意（長時間計算で有用）\\
C & $\psi\to\phi$ 写像 & artanh 前 $\psi$ クランプ \\
D & narrow-band 再初期化 & $|\phi|\le W=5\varepsilon$，pseudo-$\tau$ は $\Delta t$ と独立 \\
E & $\phi\to\psi$ 写像＋幾何（$\mathbf{n},\kappa$）& tanh 前 $\phi/2\varepsilon$ クランプ \\
F & 再初期化後 $\psi(1-\psi)$ 補正 & 必須（再初期化後必ず）\\
\hline
\end{tabular}
\end{center}

\noindent\textbf{Stage A — TVD-RK3 Shu--Osher 更新式}：
本 stage の $\psi^{n+1}$ 更新には TVD-RK3（SSPRK3）を用いる．
具体的な 3 段更新式 $\psi^{(1)},\psi^{(2)},\psi^{n+1}$ は
\S\ref{eq:cls_tvd_rk3_psi}（§7.3）に集約しており，本節では参照のみ．
時間更新作用素 $\mathcal{L}\psi = -\nabla\cdot(\mathbf{u}\psi)$ は FCCD 面値演算子による
保存形発散であり，セル境界 $x_{i+1/2}$ における面値フラックス
$(\psi u)_{f,i+1/2}$ を Hermite 補間で構成して中心発散を求める
（\S\ref{sec:advection_motivation}，§6.1）．
節点 CCD 単独（散逸なし）は $\psi$ に適用禁止
（\S\ref{sec:scheme_per_variable}，§6.0）．
% backward-compat: 旧 eq:cls_rk3_{1,2,3} は §7.3 eq:cls_tvd_rk3_psi に統合．
% 旧ラベルは外部参照保護のため alias 化：
\phantomsection\label{eq:cls_rk3_1}%
\phantomsection\label{eq:cls_rk3_2}%
\phantomsection\label{eq:cls_rk3_3}%

\noindent\textbf{Stage D — Narrow-band 再初期化}：
全域 Eikonal 反復は $\mathcal{O}(N)$ 冗長．
$|\phi|\le W=5\varepsilon$ の帯のみ
\begin{equation*}
  \partial_\tau\phi + \mathrm{sgn}(\phi_0)(|\bnabla\phi|-1) = 0
\end{equation*}
を数擬時間ステップ反復．pseudo-$\tau$ は物理 $\Delta t$ とは別管理．
詳細は \S\ref{sec:eikonal_reinit}（既存）参照．

% -------------------------------------------------------------------------------
\subsubsection{\texorpdfstring{$\psi(1-\psi)$ 質量補正（界面局在）}{ψ(1-ψ) 質量補正（界面局在）}}
\label{sec:cls_mass_correction}
% -------------------------------------------------------------------------------

各格子点での補正：
\begin{equation}
  \psi_i^{\text{new}}
  = \psi_i^{\text{old}} + \lambda_\psi\cdot\psi_i^{\text{old}}(1-\psi_i^{\text{old}})\cdot\Delta_\text{global},
  \label{eq:cls_mass_correction}
\end{equation}
ここで $\Delta_\text{global}$ は領域全体の質量偏差：
\begin{equation*}
  \Delta_\text{global}
  = \frac{\sum_i\psi_i^{\text{initial}}V_i - \sum_i\psi_i^{\text{old}}V_i}
         {\sum_i\psi_i^{\text{old}}(1-\psi_i^{\text{old}})V_i},
\end{equation*}
$\lambda_\psi\in[0.5,1.0]$．
$\psi(1-\psi)$ プロファイルは界面帯（$\psi\approx 0.5$）で最大値 0.25，
バルク（$\psi\in\{0,1\}$）で零．
したがって補正は\textbf{界面局在}し，バルクの質量保存を乱さない．

\noindent\textbf{適用タイミング}：
Stage B（移流後）と Stage F（再初期化後）の 2 回．
Stage F は必須，Stage B は省略可能だが長時間計算（$t/T_\text{char}>10$）では
drift 蓄積を防ぐため推奨．

% -------------------------------------------------------------------------------
\subsubsection{Stage A の速度場確定タイミング}
\label{sec:cls_velocity_timing}
% backward-compat: 旧 sec:cls_velocity_consistency の alias は §7 NEW (移行先) 側に置く．
% -------------------------------------------------------------------------------

Stage A の $\psi$ 移流は\textbf{直前の PPE 投影済み} $\mathbf{u}^n$ を入力として用い，
1 ステップで $\psi^{n+1}$ を確定する．ここで「PPE 投影済み速度場」を採用する根拠
（CLS と運動量予測子で同一の $\mathbf{u}^n$ を共有しないと PPE RHS 不整合を生む）と，
正しい時間ステップ ordering の図示は §~\ref{sec:cls_velocity_consistency}（§7）に集約．
本節は CLS 6 段階アルゴリズムの責務分離（mass / geometry / correction）に焦点を絞る．

% -------------------------------------------------------------------------------
\subsubsection{失敗モードと回避}
\label{sec:cls_failure_modes}
% -------------------------------------------------------------------------------

\begin{itemize}
  \item \textbf{$\psi$ に節点 CCD（散逸・面値構成なし）を適用}→急 tanh 勾配増幅→NaN．
        必ず FCCD 面値演算子（保存形フラックス；\S\ref{sec:advection_motivation}）を用いる．
  \item \textbf{再初期化前に $\psi$ から曲率計算}→
        symmetric tanh profile の曲率は最大 $\mathcal{O}(1/\varepsilon)$ で
        真の界面曲率を $\mathcal{O}(10)$ 過大評価．
        $\phi$ から $\mathbf{n}=\bnabla\phi/|\bnabla\phi|$，
        $\kappa=\bnabla\cdot\mathbf{n}$（HFE, \S\ref{sec:hfe}）を計算．
  \item \textbf{質量補正を $\phi$ に適用}→符号距離性が崩れ再初期化必須．
  \item \textbf{$\varepsilon$ の動的変更}→$\kappa$ 定義が変わり界面運動が変質．
        シミュレーション全体で固定．
\end{itemize}

\noindent 検証：バルクの $\psi\in\{0,1\}$ 位置で
Stage F 補正が恒等的にゼロになることを毎ステップ監視する
（本稿の回帰テスト）．

\medskip
\noindent\textbf{次章への接続：}
以上で CLS 6 段階アルゴリズム（Stage A--F）の責務分離（mass / geometry / correction）が確立した．
続く §~\ref{ch:per_term_spatial}（§6）では，CLS $\psi$ 移流に加え運動量 $u,v$ 移流・
粘性 $\bnabla\cdot(2\mu\bm{D})$・物性値 $\rho,\mu$・圧力 $p$ jump の各方程式項について
「場の特性に合わせて差分演算子を選択せよ」の方針で個別離散化を論じる．

% ═══════════════════════════════════════════════════════════════════════════════
